\section{A Second Theorem on the Same Foundation}
\label{sec:second}

The foundation described above was built for one theorem.  To test whether
it is a foundation rather than scaffolding, we formalized a second
finite-injury theorem against it: the Sacks Splitting Theorem
\cite{sacks1963}, in its splitting-with-non-reducibility form.

\begin{lstlisting}
theorem sacks_splitting :
    forall A : Set Nat, CE A -> not (ComputableSet A) ->
      exists A0 A1, CE A0 /\ CE A1 /\ Disjoint A0 A1 /\ A0 union A1 = A
        /\ not (A <=T A0) /\ not (A <=T A1)
\end{lstlisting}

Lowness of the halves, the full theorem's additional conclusion, requires
infinite-injury machinery and is out of scope.

The two developments are \emph{siblings}: neither imports the other, and
Lake enforces this.  Both import a shared library
\texttt{OracleComputability} holding the machine model, the use principle,
the reducibility vocabulary, the Mathlib bridge, and the monotone-stage
approximation layer.  Building the Sacks library from a clean tree compiles
no module of the Friedberg--Muchnik library.  A single root module joins
them, and contains the only declaration depending on both: the
Friedberg--Muchnik construction yields a c.e.\ non-computable set, which is
precisely the hypothesis Sacks splitting consumes, so the first theorem
certifies that the second is not vacuous.

\paragraph{What transferred.}
The entire foundation transferred unchanged.  Two items are worth naming.
The restraint lemma---a computation seen halting against a stage snapshot
survives to the limit provided nothing below its recorded use is later
enumerated---was already stated generically in the stage sequence, so the
preservation argument at the heart of Sacks reuses it verbatim.  More
interestingly, the determinism lemma \texttt{run\_halt\_unique} carries
more weight in the second theorem than in the first.  In
Friedberg--Muchnik it says only that a computation has one output.  In
Sacks it is what makes \emph{the use of the limit computation at $y$} a
well-defined number, and that number is what bounds the restraint
function.  Recording the use rather than only the output was a design
decision taken for the first theorem; it paid for a theorem it was not
designed for.

\paragraph{What did not.}
The priority layer did not transfer, and no generalization of it would
have.  The Friedberg--Muchnik injury argument counts \emph{events}: a
requirement's record carries a rank in $\{0,1,2\}$, receiving attention
raises it strictly, and only injury resets it, so once the superiors go
quiet a requirement acts at most twice more.  A splitting requirement
generates no events.  It never acts---the elements are not ours to
choose---and its restraint moves continuously with the length of
agreement, so there is nothing to count.  Its replacement is a three-part
simultaneous induction along the priority order: bounded restraints below
$j$ give finite injury to $j$; finite injury to $j$ gives satisfaction of
$j$; satisfaction of $j$ bounds $j$'s restraint.  The two arguments share a
shape---induction along the priority order, no closed-form injury
bound---and no content.  We record this as a negative result about
reuse: ``finite injury'' names a family of arguments, not a single lemma.

\paragraph{Where the hypothesis is spent.}
The one genuinely new piece of mathematics is the step at which
non-computability of $A$ is consumed.  If a requirement is never injured
after some stage and its reduction nevertheless computes $\chi_A$
everywhere, then $A$ is computable: to decide $y \in A$, search for a stage
past the last injury whose length of agreement exceeds $y$ and report
$y \in A_s$.  Such a stage exists because the reduction is total; the
answer is correct because past the last injury the agreeing computation is
frozen by the restraint lemma, so it \emph{is} the true value; and the
search is effective because the whole construction is primitive recursive.
Friedberg--Muchnik has no analogue: its requirements are satisfiable
because their witnesses are fresh, which is a property of the
construction, whereas a splitting requirement is satisfiable only because
of a property of the \emph{given} set.

\paragraph{Cost.}
The second theorem is roughly 1{,}600 lines against the first
development's 4{,}250.  The ratio is the point: the expensive
part---model, evaluator, use principle, Mathlib bridge---was already
present.  Two components of the first development turned out to be
artifacts of its own design rather than necessities.  Its ten-field state
invariant exists to keep freely chosen witnesses fresh, distinct and clear
of higher-priority restraints; a construction that does not choose its
elements needs none of it, and the two invariants that remain (the halves
stay disjoint, and their union is the current stage of $A$) are one
induction.  Its computability layer needed a shadow implementation on
plain tuples because its state types are custom structures without a
\texttt{Primcodable} instance; defining the splitting state as a tuple from
the outset removes that layer entirely.

Finally, three generic \texttt{Primrec} helpers had been declared
\texttt{private} inside the first development's computability file.
Nothing in them mentions either construction---they are what any stage
function needs to be shown computable---but \texttt{private} made them
invisible downstream.  Nothing needed generalizing; they needed exporting.
This is the kind of reusability defect that is invisible until a second
client exists.
